%% P18 --- Matrix calculus
%%
%% Stub, generated by tools/gen_stubs.py from tools/programs.json.
%% The brief below is the contract for this program: what it must argue, and
%% what it must buy the reader. Writing the program means deleting the
%% \programstub{} block and replacing it with the program -- after which this
%% file is never regenerated. If the brief turns out to be wrong once the
%% mathematics has been worked, change it in the manifest and record the
%% contradiction in CLAUDE.md under *Resolved questions*.

\program{Matrix calculus}
\label{prog:P18}

\programstub{%
Argument: differentiating with respect to a vector or a matrix is
bookkeeping plus a layout convention, and most of the pain in the literature
is the convention rather than the mathematics. Payoff: the identities the
reader actually needs, derived once each: the gradient of \code{Wx} with
respect to \code{W}; of a squared error; of a softmax; of a log-softmax; of
layer normalisation. The headline: **the gradient of cross-entropy through a
softmax is \code{p - y}**, the single most reused fact in applied machine
learning, derived in full \dash{} and the reason the two operations are
fused in every serious implementation. Also a versionbox on numerator versus
denominator layout, because the reader will meet both in the same week.

FORWARD REFERENCE, declared: the softmax--cross-entropy gradient is this
program's most reused derivation, and cross-entropy is not defined until
P30. Give it a definitional frame here -- cross-entropy as minus the sum of
y log p, stated as a definition to be justified later -- and have P26 and
P30 each carry a frame returning to it.

\medskip
\textbf{Planned length:} 60 frames.
}
